% Repository management, lecture 1 of 2 -- the repository as the centre of
% the project: philosophy, Git in broad strokes, organisation, protection.
%
% Companion to ../L01-repository.md, which is the source of truth for structure
% and theses. Slide text and speaker notes are in Russian
% because that is the content; everything else in this file is English.
%
% Build:  make          (latexmk -pdf; engine: pdflatex)
%         make notes    (same deck plus speaker notes on a second screen)
%
% The history on the diagrams is real: every hash comes from git 2.43.0 run
% on a throwaway HelloWorld repository (merge, squash and rebase included).
%
% House rules: ../CLAUDE.md and ../../../../CLAUDE.md.
% Theme: ../../../../../dotnet-course/theme/.

\documentclass[aspectratio=169,10pt,t]{beamer}

\usetheme{dotnet}
\usetikzlibrary{arrows.meta,positioning,calc}

\dotnetlecture{Л1}

\ifdefined\shownotes
  \setbeameroption{show notes on second screen=right}
\fi

\title[Л1. Репозиторий]{Управление репозиторием: центр мироздания проекта}
\author{Дмитрий Курочкин}
\date{}

\tikzset{
  box/.style={draw=hairline,line width=0.8pt,fill=codebg,
              minimum height=0.95cm,align=center,inner sep=0.4em},
  ar/.style={-{Straight Barb[length=1.8mm]},ink,line width=0.7pt},
  % vertical history diagrams: a commit is a dot on its branch's line
  dot/.style={circle,draw=accent,line width=1.0pt,fill=paper,
              minimum size=0.36cm,inner sep=0pt},
  bdot/.style={dot,draw=alt},
  gdot/.style={dot,draw=hairline,dashed},
  ml/.style={accent,line width=1.6pt},
  bl/.style={alt,line width=1.6pt},
  gl/.style={hairline,line width=1.2pt,dashed},
  cl/.style={font=\scriptsize,align=left,anchor=west},
  cr/.style={font=\scriptsize,align=right,anchor=east},
  ref/.style={font=\scriptsize\bfseries},
}

% Fit inside a 2.6cm square rather than scaling by width alone: the supplied
% codes are not all square, and width-only scaling makes a taller one overrun
% the row.
\newcommand{\qrbox}[1]{%
  \IfFileExists{img/#1.png}%
    {\includegraphics[width=2.6cm,height=2.6cm,keepaspectratio]{img/#1.png}}%
    {\fcolorbox{hairline}{paper}{\parbox[c][2.5cm][c]{2.5cm}%
       {\centering\color{muted}\scriptsize QR}}}}

\begin{document}

% ------------------------------------------------------------------ title --
\begin{frame}[plain]
  \titlepagednc{Управление репозиторием · лекция 1 из 2}{Репозиторий ---
    центр мироздания проекта}{Философия, Git, организация, защита}{Дмитрий
    Курочкин}
\end{frame}
\note{80 минут. Это не курс <<как пользоваться Git>>. Git сегодня будет, но
  в общих чертах: команды читаются за вечер, а мышление -- нет.}

% ================================================= 0. WHY ==================
\section{0. Зачем это всё}

\begin{frame}{Контакты}
  \slidetop
  \begin{center}
    \begin{tabular}{@{}c@{\hspace{1.6cm}}c@{\hspace{1.6cm}}c@{}}
      \qrbox{github-qr} & \qrbox{telegram-qr} & \qrbox{nuget-qr} \\[0.8em]
      {\color{accent}\bfseries GitHub} &
      {\color{accent}\bfseries Telegram} &
      {\color{accent}\bfseries NuGet} \\
    \end{tabular}
  \end{center}

  \slidegap
  \begin{center}
    {\large\texttt{dmitry.kurochkin0@gmail.com}}
  \end{center}
\end{frame}
\note{Все ссылки здесь, на одном слайде. Регалии -- голосом, одной фразой.
  Дальше сразу вопрос в зал.}

\begin{frame}[plain]
  \vote{Кто собирается быть разработчиком ПО?\\[0.4em] А кто ---
    вайбкодером?}
\end{frame}
\note{Два голосования подряд. Не осуждать вторых: ИИ -- это навигатор, без
  него уже странно. Вопрос в том, кто отвечает, когда навигатор ошибся.}

\begin{frame}{Код подешевел}
  \slidetop
  \claim{Написать код больше не трудно. Трудно --- \hi{не пустить плохой код
    в продукт}.}

  \versus{что осталось дорогим}{%
    \begin{cbul}{accent}
      \item проверить, что это работает
      \item принять или отклонить
      \item не сломать тех, кто уже пользуется
      \item отвечать за результат
    \end{cbul}}{что стало дешёвым}{%
    \begin{cbul}{alt}
      \item написать функцию
      \item предложить изменение
      \item сгенерировать тест
      \item придумать API или зависимость
    \end{cbul}}
\end{frame}
\note{Изменение сейчас ничего не стоит: агент предложит десять за час. Значит,
  ценность переехала от написания к контролю.
  Вся лекция про левую колонку.}

\begin{frame}[plain]
  \statement{Профессионализм --- это не размер кода, который вы способны
    написать, а \hi{размер хаоса}, который вы способны контролировать.}
\end{frame}
\note{Главная мысль обеих лекций. Представьте проект: 50 разработчиков и
  агентов. Кто-то ошибается, кто-то сдаёт плохой код, кто-то забывает
  документацию, ИИ придумывает несуществующий пакет. Нужен механизм, который
  не зависит от доверия к конкретному человеку.}

\begin{frame}{Пять навыков}
  \slidetop
  \claim{Две лекции --- один маршрут: \hi{от хаоса к системе}.}

  \begin{center}
    \begin{tikzpicture}[font=\small,
      sk/.style={box,text width=2.05cm,minimum height=1.25cm}]
      \node[sk,draw=accent,line width=1.0pt] (a) {\hi{философия}};
      \node[sk,draw=accent,line width=1.0pt,right=0.45cm of a] (b)
        {\hi{организация}};
      \node[sk,draw=accent,line width=1.0pt,right=0.45cm of b] (c)
        {\hi{защита}};
      \node[sk,draw=alt,line width=1.0pt,right=0.45cm of c] (d)
        {\lo{управление задачами}};
      \node[sk,draw=alt,line width=1.0pt,right=0.45cm of d] (e)
        {\lo{поставка}};
      \draw[ar] (a) -- (b); \draw[ar] (b) -- (c);
      \draw[ar] (c) -- (d); \draw[ar] (d) -- (e);
      \node[accent,font=\scriptsize\bfseries,below=0.7em] at (b.south)
        {сегодня};
      \node[alt,font=\scriptsize\bfseries,below=0.7em]
        at ($(d.south)!0.5!(e.south)$) {следующая лекция};
    \end{tikzpicture}
  \end{center}

  \slidegap
  \aside{В конце вы должны уметь взять проект, в котором царит хаос, и
    превратить его в систему, где изменения контролируются, воспроизводятся
    и проверяются.}
\end{frame}
\note{Сегодня: что такое репозиторий, сам Git, как репозиторий устроен и как
  он защищается. В следующий раз: лидер, задачи, деньги и поставка.}

% ================================================= 1. PHILOSOPHY ===========
\section{1. Философия}

\begin{frame}[plain]
  \sectionopener{1}{Философия}{Что должно лежать в репозитории, кроме
    \hi{кода}?}
\end{frame}
\note{Бюджет: 15 минут. Вопрос в зал, собрать ответы голосом.}

\begin{frame}{Репозиторий --- центр мироздания проекта}
  \slidetop
  \claim{Всё, что нужно знать о продукте, лежит \hi{в одном месте} --- или
    не существует.}

  \begin{center}
    \begin{tikzpicture}[font=\small,
      sat/.style={box,minimum width=2.35cm,minimum height=0.8cm}]
      \node[box,draw=accent,line width=1.2pt,minimum width=3.0cm,
            minimum height=1.1cm] (r) {\textbf{\hi{репозиторий}}};
      \node[sat] (s1) at (-5.0, 1.15) {код};
      \node[sat] (s2) at (-5.0, 0.00) {тесты};
      \node[sat] (s3) at (-5.0,-1.15) {документация};
      \node[sat] (s4) at ( 0.0, 1.50) {задачи};
      \node[sat] (s5) at ( 0.0,-1.50) {решения};
      \node[sat] (s6) at ( 5.0, 1.15) {CI и правила};
      \node[sat] (s7) at ( 5.0, 0.00) {зависимости};
      \node[sat] (s8) at ( 5.0,-1.15) {релизы};
      \foreach \s in {s1,s2,s3,s4,s5,s6,s7,s8}
        \draw[hairline,line width=0.8pt] (r) -- (\s);
    \end{tikzpicture}
  \end{center}
\end{frame}
\note{Не только исходники. Задачи, документация, конфигурация, lockfile,
  тесты, CI, правила участия, владельцы кода, релизы, changelog и полный
  след изменений. Репозиторий отвечает на вопросы: что это, как войти, как
  собрать, как проверить, как изменить, как доставить.}

\begin{frame}{Проверка: человек с улицы}
  \slidetop
  \claim{Новый участник \hi{ничего не спрашивает}. Он клонирует --- и
    работает.}

  \begin{tabular}{@{}p{4.6cm}p{8.6cm}@{}}
    \textbf{\hi{что это?}}          & \texttt{README} \\[0.5em]
    \textbf{\hi{как запустить?}}    & одна команда, одинаковая у всех \\[0.5em]
    \textbf{\hi{что делать?}}       & открытые задачи \\[0.5em]
    \textbf{\hi{как сдать работу?}} & \texttt{CONTRIBUTING} --- правила участия \\[0.5em]
    \textbf{\hi{почему так сделано?}} & история и записи о решениях \\
  \end{tabular}

  \slidegap
  \aside{Этот человек может быть и \lo{агентом}. Ему тем более некого
    спросить в коридоре.}
\end{frame}
\note{Self-contained: репозиторий не требует устного предания. Если что-то
  живёт в голове у одного человека -- этого нет.}

\begin{frame}{Работа --- это поставка артефактов}
  \slidetop
  \claim{<<Делал, думал, анализировал>> --- не работа, пока от неё нет
    \hi{цифрового следа}.}

  \begin{tabular}{@{}p{4.6cm}p{8.6cm}@{}}
    \textbf{\lo{<<изучал>>}}        & \hi{документ} или запись о решении \\[0.5em]
    \textbf{\lo{<<нашёл проблему>>}} & \hi{задача} в репозитории \\[0.5em]
    \textbf{\lo{<<придумал>>}}      & \hi{черновик изменения}, который видно всем \\[0.5em]
    \textbf{\lo{<<обсудили>>}}      & \hi{комментарий} в задаче \\[0.5em]
    \textbf{\lo{<<почти готово>>}}  & \hi{коммит}, который можно
      посмотреть \\
  \end{tabular}
\end{frame}
\note{Вся работа -- поставка цифровых артефактов, и ничего другого нет.
  Остальное -- имитация. Изучал? Покажи артефакт. Это не про недоверие, это
  про то, что результат должен пережить вашу память.}

\begin{frame}{Жалоба --- это вклад}
  \slidetop
  \claim{<<Я не могу запустить проект>>.}

  \versus{задача в репозитории}{%
    \begin{cbul}{accent}
      \item видят все, включая будущих
      \item у проблемы есть номер и ссылка
      \item чинится причина: \hi{плохая документация}
      \item автору говорят спасибо
    \end{cbul}}{личное сообщение}{%
    \begin{cbul}{alt}
      \item видит один человек
      \item ответ исчезает в переписке
      \item следующий новичок \lo{спросит то же}
      \item следа нет
    \end{cbul}}
\end{frame}
\note{Человек пришёл и не смог запустить -- он не пишет в чат, он оформляет
  задачу в репозитории. Это вклад: у нас плохие доки, он нам помог. Жалуйтесь постоянно --
  но в форме артефакта. Личное общение создаёт почву для отсутствия следа.}

\begin{frame}[plain]
  \keyline{Чего нет в репозитории --- того \hi{не было}.}
\end{frame}
\note{Чтобы так жить, нужен инструмент, который помнит всё. Он называется
  Git.}

% ================================================= 2. GIT ==================
\section{2. Git в общих чертах}

\begin{frame}[plain]
  \sectionopener{2}{Git в общих чертах}{Команды вы выучите за вечер.
    Сейчас --- \hi{модель}, без которой команды не имеют смысла.}
\end{frame}
\note{Бюджет: 24 минуты. Кто уже пользовался Git? Кто делал git push --force
  и жалел?}

\begin{frame}{История --- это цепочка снимков}
  \slidetop
  \claim{Коммит --- \hi{снимок всего проекта} и ссылка на предыдущий
    снимок.}

  \begin{center}
    \begin{tikzpicture}[font=\small,
      cmt/.style={box,minimum width=2.3cm,minimum height=1.15cm,
                 font=\footnotesize}]
      \node[cmt] (c1) {\texttt{cb5849b}\\print greeting};
      \node[cmt,right=1.0cm of c1] (c2) {\texttt{408319b}\\add readme};
      \node[cmt,right=1.0cm of c2] (c3) {\texttt{983e9ae}\\describe usage};
      \draw[ar] (c2) -- (c1); \draw[ar] (c3) -- (c2);
      \node[accent,font=\scriptsize\bfseries,above=0.9em] (m) at (c3.north)
        {main};
      \draw[accent,line width=0.7pt,-{Straight Barb[length=1.4mm]}]
        (m) -- (c3.north);
      \node[muted,font=\scriptsize,below=0.5em] at (c2.south)
        {стрелка --- <<мой родитель>>};
    \end{tikzpicture}
  \end{center}

  \slidegap
  \aside{Имя коммита --- хеш его содержимого. Изменить прошлое незаметно
    нельзя: \hi{изменится имя} и всё, что после.}
\end{frame}
\note{Не <<разница между файлами>>, а снимок. Отсюда всё остальное: ветка --
  просто указатель на коммит, история -- граф, а не папка с копиями
  final-final-2.}

\begin{frame}{Цепочка коммитов --- это ветка}
  \slidetop
  \begin{minipage}[t]{0.40\textwidth}
    \vspace{0pt}
    \centering
    \begin{tikzpicture}[font=\footnotesize,
      dot/.style={circle,draw=accent,line width=1.0pt,fill=paper,
                  minimum size=0.42cm,inner sep=0pt}]
      \draw[accent,line width=1.6pt] (0,-0.5) -- (0,4.2);
      \node[dot] (c1) at (0,0.2) {};
      \node[dot] (c2) at (0,1.6) {};
      \node[dot,fill=accent] (c3) at (0,3.0) {};
      \foreach \c/\h/\m in {c1/cb5849b/print greeting,
                            c2/408319b/add readme,
                            c3/983e9ae/describe usage} {
        \node[anchor=west,align=left] at ([xshift=0.35cm]\c.east)
          {\texttt{\h}\\{\color{muted}\m}};
      }
      \node[accent,font=\scriptsize\bfseries,anchor=east] (m)
        at ([xshift=-0.55cm]c3.west) {main};
      \draw[accent,line width=0.7pt,-{Straight Barb[length=1.4mm]}]
        (m) -- (c3);
      \draw[muted,line width=0.6pt,-{Straight Barb[length=1.6mm]}]
        (-1.9,0.0) -- node[above,sloped,font=\scriptsize,text=muted]
        {время} (-1.9,2.2);
    \end{tikzpicture}
  \end{minipage}\hfill
  \begin{minipage}[t]{0.55\textwidth}
    \vspace{0.6em}
    \claim{\raggedright Каждый коммит знает предыдущий. Идите по этим
      ссылкам --- получится \hi{ветка}.}

    \begin{cbul}{accent}
      \item новый коммит встаёт \hi{сверху}
      \item имя ветки указывает на \hi{последний} коммит
      \item всё, что ниже, --- её история
    \end{cbul}

    \vspace{0.9em}
    \aside{Так же рисует историю сам Git: новое сверху, старое снизу.}
  \end{minipage}
\end{frame}
\note{Первая вертикальная картинка -- дальше в терминале история будет
  выглядеть так же. main -- не папка и не копия проекта, а имя для верхнего
  коммита; вся линия под ним получается из ссылок на родителя.}

\begin{frame}{Четыре места, где живёт ваш код}
  \slidetop
  \begin{center}
    \begin{tikzpicture}[font=\small,
      pl/.style={box,minimum width=2.45cm,minimum height=1.2cm}]
      \node[pl] (w) {рабочая\\папка};
      \node[pl,right=1.2cm of w] (i) {индекс};
      \node[pl,right=1.2cm of i] (l) {локальный\\репозиторий};
      \node[pl,right=1.2cm of l,draw=accent,line width=1.0pt] (r)
        {\hi{удалённый}\\\hi{репозиторий}};
      \draw[ar] (w) -- node[above,font=\scriptsize\ttfamily] {add} (i);
      \draw[ar] (i) -- node[above,font=\scriptsize\ttfamily] {commit} (l);
      \draw[ar] (l) -- node[above,font=\scriptsize\ttfamily] {push} (r);
      \draw[ar] (r.south) -- ++(0,-0.75) -|
        node[pos=0.25,below,font=\scriptsize\ttfamily] {pull} (w.south);
    \end{tikzpicture}
  \end{center}

  \slidegap
  \begin{tabular}{@{}p{3.6cm}p{9.6cm}@{}}
    \textbf{\hi{индекс}} & что именно войдёт в следующий коммит --- вы
      выбираете сами \\[0.4em]
    \textbf{\hi{локальный}} & вся история у вас на диске, работает без
      сети \\[0.4em]
    \textbf{\hi{удалённый}} & общий для команды; правда проекта --- там \\
  \end{tabular}
\end{frame}
\note{Индекс -- то, что отличает Git от <<сохранить всё>>: в коммит идёт не
  всё, что вы наменяли, а то, что вы осознанно выбрали. Это пригодится через
  два слайда.}

\begin{frame}{Первый коммит}
  \slidetop
  \claim{Четыре команды --- и снимок \hi{записан в историю}.}

  \begin{tabular}{@{}p{8.1cm}p{4.9cm}@{}}
    \texttt{git init}                 & создать репозиторий в папке \\[0.7em]
    \texttt{git status}               & что изменилось и что выбрано \\[0.7em]
    \texttt{git add Program.cs}       & \hi{выбрать} файл для коммита \\[0.7em]
    \texttt{git commit -m \textquotedbl feat: print greeting\textquotedbl}
                                      & \hi{записать} снимок с сообщением \\
  \end{tabular}

  \slidegap
  \aside{В коммит попадает только выбранное. В папке два файла, в снимке
    \hi{один}: \texttt{README.md} ждёт своего коммита и своего сообщения.}
\end{frame}
\note{Проговорить по строке и показать вживую в терминале. git status --
  команда, которую набирают чаще всех остальных: она ничего не меняет и
  всегда говорит, где вы.}

\begin{frame}{Один коммит --- одна логическая мысль}
  \slidetop
  \claim{Каждый коммит \hi{собирается, проходит тесты} и объясняет себя
    сам.}

  \versus{так}{%
    \begin{cbul}{accent}
      \item \texttt{feat: greet by name}
      \item \texttt{fix: handle empty name}
      \item \texttt{docs: describe usage}
      \item \texttt{test: cover empty name}
    \end{cbul}}{не так}{%
    \begin{cbul}{alt}
      \item \texttt{fix}
      \item \texttt{update}
      \item \texttt{final version 2}
      \item \texttt{new feature + fixes + refactoring}
    \end{cbul}}

  \slidegap
  \aside{Сообщение отвечает на вопрос \hi{<<что изменилось и зачем>>}.
    Его будут читать через год --- скорее всего, вы.}
\end{frame}
\note{Атомарный коммит. Формат слева -- Conventional Commits: тип, двоеточие,
  что сделано. Из таких сообщений потом автоматически собирается changelog и
  считается номер версии -- об этом в следующей лекции.}

\begin{frame}{Ветка под задачу}
  \slidetop
  \begin{minipage}[t]{0.47\textwidth}
    \vspace{0pt}
    \centering
    \begin{tikzpicture}
      \draw[ml] (0,-0.4) -- (0,3.4);
      \draw[bl] (0,0.9) to[out=70,in=-90] (1.5,1.9) -- (1.5,3.4);
      \node[dot] (c1) at (0,0.0) {};
      \node[dot] (c2) at (0,0.9) {};
      \node[dot,fill=accent] (c3) at (0,2.7) {};
      \node[bdot] (b1) at (1.5,1.9) {};
      \node[bdot,fill=alt] (b2) at (1.5,2.8) {};
      \node[cr] at ([xshift=-0.25cm]c1.west) {\texttt{cb5849b}\\{\color{muted}print greeting}};
      \node[cr] at ([xshift=-0.25cm]c2.west) {\texttt{408319b}\\{\color{muted}add readme}};
      \node[cr] at ([xshift=-0.25cm]c3.west) {\texttt{983e9ae}\\{\color{muted}describe usage}};
      \node[cl] at ([xshift=0.25cm]b1.east) {\texttt{3508968}\\{\color{muted}greet by name}};
      \node[cl] at ([xshift=0.25cm]b2.east) {\texttt{c320998}\\{\color{muted}return exit code}};
      \node[ref,accent,anchor=south] at (0,3.45) {main};
      \node[ref,alt,anchor=south west] at (1.2,3.45) {feat/greet-by-name};
    \end{tikzpicture}
  \end{minipage}\hfill
  \begin{minipage}[t]{0.49\textwidth}
    \vspace{0.6em}
    \claim{\raggedright Новая ветка --- \lo{вторая линия} от того же
      коммита.}

    \texttt{git switch -c feat/greet-by-name}

    \vspace{0.9em}
    \begin{cbul}{accent}
      \item ветка не копирует проект: это ещё одно \hi{имя}
      \item ваши коммиты ложатся на \lo{свою линию}
      \item main тем временем \hi{уходит вперёд}
    \end{cbul}
  \end{minipage}
\end{frame}
\note{Ветка ничего не стоит, поэтому на каждую задачу -- своя. Пока вы
  работали, в main появился 983e9ae: это нормальное состояние, и чем дольше
  живёт ветка, тем дальше уходит main. Две линии надо когда-то соединить --
  способов три.}

\begin{frame}{Слияние: merge commit}
  \slidetop
  \begin{minipage}[t]{0.47\textwidth}
    \vspace{0pt}
    \centering
    \begin{tikzpicture}
      \draw[ml] (0,-0.4) -- (0,4.0);
      \draw[bl] (0,0.9) to[out=70,in=-90] (1.5,1.7) -- (1.5,2.6)
        to[out=90,in=-70] (0,3.5);
      \node[dot] (c1) at (0,0.0) {};
      \node[dot] (c2) at (0,0.9) {};
      \node[dot] (c3) at (0,2.5) {};
      \node[cr] at ([xshift=-0.25cm]c1.west) {\texttt{cb5849b}\\{\color{muted}print greeting}};
      \node[cr] at ([xshift=-0.25cm]c2.west) {\texttt{408319b}\\{\color{muted}add readme}};
      \node[cr] at ([xshift=-0.25cm]c3.west) {\texttt{983e9ae}\\{\color{muted}describe usage}};
      \node[bdot] (b1) at (1.5,1.7) {};
      \node[bdot] (b2) at (1.5,2.6) {};
      \node[dot,fill=accent] (m) at (0,3.5) {};
      \node[cl] at ([xshift=0.25cm]b1.east) {\texttt{3508968}\\{\color{muted}greet by name}};
      \node[cl] at ([xshift=0.25cm]b2.east) {\texttt{c320998}\\{\color{muted}return exit code}};
      \node[cr] at ([xshift=-0.25cm]m.west) {\texttt{845636a}\\{\color{accent}merge commit}};
      \node[ref,accent,anchor=south] at (0,4.05) {main};
    \end{tikzpicture}
  \end{minipage}\hfill
  \begin{minipage}[t]{0.49\textwidth}
    \vspace{0.6em}
    \claim{\raggedright Две линии сходятся в \hi{одном коммите} с двумя родителями.}

    \texttt{git merge feat/greet-by-name}

    \vspace{0.9em}
    \begin{cbul}{accent}
      \item история сохранена \hi{как была}
      \item видно, где ветка началась и где влилась
      \item цена: граф ветвится
    \end{cbul}
  \end{minipage}
\end{frame}
\note{Первый способ. Ничего не теряется: оба коммита ветки остаются в
  истории, а над ними появляется коммит слияния. На большом проекте с сотней
  веток граф становится трудно читать.}

\begin{frame}{Слияние: squash}
  \slidetop
  \begin{minipage}[t]{0.42\textwidth}
    \vspace{0pt}
    \centering
    \begin{tikzpicture}
      \draw[ml] (0,-0.4) -- (0,4.0);
      \draw[gl] (0,0.9) to[out=70,in=-90] (1.5,1.7) -- (1.5,2.9);
      \node[dot] (c1) at (0,0.0) {};
      \node[dot] (c2) at (0,0.9) {};
      \node[dot] (c3) at (0,2.5) {};
      \node[cr] at ([xshift=-0.25cm]c1.west) {\texttt{cb5849b}\\{\color{muted}print greeting}};
      \node[cr] at ([xshift=-0.25cm]c2.west) {\texttt{408319b}\\{\color{muted}add readme}};
      \node[cr] at ([xshift=-0.25cm]c3.west) {\texttt{983e9ae}\\{\color{muted}describe usage}};
      \node[gdot] (b1) at (1.5,1.7) {};
      \node[gdot] (b2) at (1.5,2.6) {};
      \node[dot,fill=accent] (m) at (0,3.5) {};
      \node[cl,text=muted] at ([xshift=0.25cm]b1.east) {\texttt{3508968}};
      \node[cl,text=muted] at ([xshift=0.25cm]b2.east) {\texttt{c320998}};
      \draw[alt,line width=0.7pt,dashed,-{Straight Barb[length=1.6mm]}]
        (1.5,2.95) to[out=90,in=0] (m.east);
      \node[cl,text=alt] at (1.58,3.3) {изменения\\обоих коммитов};
      \node[cr] at ([xshift=-0.25cm]m.west) {\texttt{5e77d87}\\{\color{accent}greet by name}};
      \node[ref,accent,anchor=south] at (0,4.05) {main};
    \end{tikzpicture}
  \end{minipage}\hfill
  \begin{minipage}[t]{0.56\textwidth}
    \vspace{0.6em}
    \claim{\raggedright Вся ветка превращается в \hi{один коммит} в main.}

    \texttt{git merge -{}-squash feat/greet-by-name}

    \vspace{0.9em}
    \begin{cbul}{accent}
      \item main \hi{прямой}: одна задача --- один коммит
      \item черновые коммиты ветки не засоряют историю
      \item цена: шаги внутри ветки \lo{теряются}
    \end{cbul}
  \end{minipage}
\end{frame}
\note{Второй способ. Коммиты ветки в main не попадают: попадает один новый,
  в котором лежат изменения их обоих. Хорош, когда в ветке <<fix>>, <<ещё
  fix>>, <<теперь точно>>.}

\begin{frame}{Слияние: rebase}
  \slidetop
  \begin{minipage}[t]{0.47\textwidth}
    \vspace{0pt}
    \centering
    \begin{tikzpicture}
      \draw[ml] (0,-0.4) -- (0,4.5);
      \draw[gl] (0,0.9) to[out=70,in=-90] (1.5,1.7) -- (1.5,2.3);
      \node[dot] (c1) at (0,0.0) {};
      \node[dot] (c2) at (0,0.9) {};
      \node[dot] (c3) at (0,2.2) {};
      \node[cr] at ([xshift=-0.25cm]c1.west) {\texttt{cb5849b}\\{\color{muted}print greeting}};
      \node[cr] at ([xshift=-0.25cm]c2.west) {\texttt{408319b}\\{\color{muted}add readme}};
      \node[cr] at ([xshift=-0.25cm]c3.west) {\texttt{983e9ae}\\{\color{muted}describe usage}};
      \node[gdot] (b1) at (1.5,1.5) {};
      \node[gdot] (b2) at (1.5,2.1) {};
      \node[ref,accent,anchor=west] at (0.28,2.5) {main};
      \node[bdot] (r1) at (0,3.1) {};
      \node[bdot,fill=alt] (r2) at (0,4.0) {};
      \node[cl,text=muted] at ([xshift=0.2cm]b1.east) {\texttt{3508968}};
      \node[cl,text=muted] at ([xshift=0.2cm]b2.east) {\texttt{c320998}};
      \draw[alt,line width=0.7pt,dashed,-{Straight Barb[length=1.6mm]}]
        (1.5,2.35) to[out=90,in=0] (r1.east);
      \node[cl,text=alt] at (0.9,3.6) {те же изменения,\\новые коммиты};
      \node[cr] at ([xshift=-0.25cm]r1.west) {\texttt{68fa711}\\{\color{muted}greet by name}};
      \node[cr] at ([xshift=-0.25cm]r2.west) {\texttt{bffa5e2}\\{\color{muted}return exit code}};
      \node[ref,alt,anchor=south] at (0,4.55) {feat/greet-by-name};
    \end{tikzpicture}
  \end{minipage}\hfill
  \begin{minipage}[t]{0.49\textwidth}
    \vspace{0.6em}
    \claim{\raggedright Коммиты ветки \lo{переставляются} поверх свежего main.}

    \texttt{git rebase main}

    \vspace{0.9em}
    \begin{cbul}{accent}
      \item история \hi{прямая}, и каждый шаг сохранён
      \item коммиты получают \lo{новые имена}: это другие коммиты
      \item нельзя делать с веткой, которой уже пользуются другие
    \end{cbul}
  \end{minipage}
\end{frame}
\note{Третий способ. Git берёт изменения каждого коммита ветки и записывает
  их заново, начиная от верхнего коммита main. Хеши другие, потому что
  другой родитель. После этого main просто передвигается вверх -- без
  коммита слияния. Договоритесь о способе один раз на проект.}

\begin{frame}{Команды на каждый день}
  \slidetop
  {\small
  \begin{tabular}{@{}p{4.9cm}p{8.3cm}@{}}
    \texttt{git clone <url>}          & получить репозиторий целиком, с
      историей \\[0.38em]
    \texttt{git switch -c <branch>}   & новая ветка под задачу \\[0.38em]
    \texttt{git status}, \texttt{git diff} & что я изменил \\[0.38em]
    \texttt{git add}, \texttt{git commit}  & выбрать и записать снимок \\[0.38em]
    \texttt{git push}                 & отправить ветку на сервер \\[0.38em]
    \texttt{git pull}                 & забрать чужие изменения \\[0.38em]
    \texttt{git log --oneline --graph} & прочитать историю \\[0.38em]
    \texttt{git revert <commit>}      & \hi{отменить коммит новым
      коммитом} \\[0.38em]
    \texttt{git push --force} в общую ветку & \lo{переписать чужую историю
      --- нельзя} \\
  \end{tabular}}
\end{frame}
\note{Девять строк покрывают девяносто процентов дней. Revert, а не
  переписывание: в общей истории ошибки исправляют вперёд. Force push в
  общую ветку -- единственная команда на слайде, отмеченная кирпичным.}

\begin{frame}{Чему в репозитории не место}
  \slidetop
  \claim{Git помнит всё \hi{навсегда}. Поэтому не всё в него кладут.}

  \begin{tabular}{@{}p{4.6cm}p{8.6cm}@{}}
    \textbf{\lo{результаты сборки}} & \texttt{bin}, \texttt{obj} ---
      в \texttt{.gitignore}: они получаются из кода \\[0.55em]
    \textbf{\lo{пароли и ключи}}    & никогда; попавший в историю секрет
      считается украденным \\[0.55em]
    \textbf{\lo{большие бинарные файлы}} & \hi{Git LFS}: в истории
      лежит указатель, сам файл --- в отдельном хранилище \\
  \end{tabular}

  \slidegap
  \aside{Удалить файл следующим коммитом --- не значит удалить его из
    истории.}
\end{frame}
\note{Git хранит снимки, а бинарный файл не сжимается разницей: десять
  версий видео -- десять полных копий у каждого, кто клонирует. Секрет:
  не удалять коммитом, а отзывать и перевыпускать ключ.}

% ================================================= 3. ORGANISATION =========
\section{3. Организация репозитория и GitHub}

\begin{frame}[plain]
  \sectionopener{3}{Организация репозитория и GitHub}{Одному хватает Git.
    Когда разработчиков много, нужно \hi{общее место} и \hi{правила}.}
\end{frame}
\note{Бюджет: 22 минуты. Сначала -- где живёт общий репозиторий и что умеет
  площадка, потом -- что в этом репозитории должно лежать.}

\begin{frame}{Когда разработчиков много}
  \slidetop
  \claim{У каждого --- \hi{полная копия} с историей. Общая правда ---
    \hi{одна}, на сервере.}

  \begin{center}
    \begin{tikzpicture}[font=\small,
      dv/.style={box,minimum width=2.9cm,minimum height=1.0cm}]
      \node[dv,draw=accent,line width=1.0pt,minimum width=3.6cm] (o)
        at (0,2.1) {\hi{удалённый репозиторий}\\{\scriptsize\color{muted}origin}};
      \node[dv] (a) at (-4.3,0) {ваш ноутбук};
      \node[dv] (b) at ( 0.0,0) {коллега};
      \node[dv] (c) at ( 4.3,0) {агент};
      \draw[ar] (a.north) -- node[left=0.3em,font=\scriptsize\ttfamily]
        {push} (o.south west);
      \draw[ar] (o.south) -- node[right,font=\scriptsize\ttfamily]
        {pull} (b.north);
      \draw[ar] (o.south east) -- node[right=0.3em,font=\scriptsize\ttfamily]
        {clone} (c.north);
    \end{tikzpicture}
  \end{center}

  \slidegap
  \aside{Никто не пересылает файлы друг другу. Все отправляют коммиты
    \hi{в одно место} и забирают чужие оттуда же.}
\end{frame}
\note{Clone -- получить копию целиком. Push -- отправить свои коммиты.
  Pull -- забрать чужие. Локальная копия у каждого полная, поэтому работать
  можно без сети, а сервер -- это место встречи, а не единственное
  хранилище.}

\begin{frame}{GitHub --- это не Git}
  \slidetop
  \versus{Git}{%
    \begin{cbul}{accent}
      \item программа у вас на машине
      \item коммиты, ветки, слияние
      \item работает без сети
      \item про людей \hi{ничего не знает}
    \end{cbul}}{GitHub}{%
    \begin{cbul}{alt}
      \item площадка в сети
      \item хранит удалённый репозиторий
      \item задачи, pull request, проверки
      \item знает, \lo{кому что можно}
    \end{cbul}}

  \slidegap
  \aside{Git отвечает за историю. Всё, что касается \hi{совместной работы},
    --- это площадка: GitHub, GitLab и другие.}
\end{frame}
\note{Частая путаница первого курса. Pull request, issue, права доступа --
  этого в Git нет. Мы говорим про GitHub, потому что на нём живёт
  большинство открытых проектов; идеи те же везде.}

\begin{frame}{Issue}
  \slidetop
  \claim{Любая работа начинается с \hi{записи}: что не так или что нужно
    сделать.}

  \begin{tabular}{@{}p{4.6cm}p{8.6cm}@{}}
    \textbf{\hi{номер}}       & \texttt{\#42} --- на него можно сослаться откуда
      угодно \\[0.5em]
    \textbf{\hi{обсуждение}}  & вопросы и ответы остаются под задачей,
      а не в переписке \\[0.5em]
    \textbf{\hi{метки}}       & ошибка, улучшение, документация \\[0.5em]
    \textbf{\hi{исполнитель}} & видно, кто взял и взял ли кто-нибудь \\
  \end{tabular}

  \slidegap
  \aside{Issue может завести кто угодно: разработчик, пользователь,
    новичок, который не смог запустить проект.}
\end{frame}
\note{Коротко: это та самая жалоба в форме артефакта из первого раздела.
  Как ставить задачи и как ими управлять -- следующая лекция.}

\begin{frame}{Pull request}
  \slidetop
  \claim{В main никто не пишет напрямую. Туда попадают через \hi{pull
    request}.}

  \begin{center}
    \begin{tikzpicture}[font=\small,
      st/.style={box,minimum width=1.75cm,minimum height=1.1cm}]
      \node[st] (a) {задача};
      \node[st,right=0.55cm of a] (b) {ветка};
      \node[st,right=0.55cm of b] (c) {коммиты};
      \node[st,right=0.55cm of c,draw=accent,line width=1.0pt] (d)
        {\hi{pull}\\\hi{request}};
      \node[st,right=0.55cm of d] (e) {проверки\\и review};
      \node[st,right=0.55cm of e] (f) {main};
      \foreach \x/\y in {a/b,b/c,c/d,d/e,e/f} \draw[ar] (\x) -- (\y);
      \draw[alt,line width=0.7pt,-{Straight Barb[length=1.8mm]}]
        (e.south) -- ++(0,-0.6) -|
        node[pos=0.25,below,font=\scriptsize,text=alt] {не прошло ---
        исправить} (c.south);
    \end{tikzpicture}
  \end{center}

  \slidegap
  \aside{Pull request --- это \hi{просьба}: <<возьмите мою ветку в main>>.
    Просьбу можно отклонить.}
\end{frame}
\note{PR -- не команда Git, а механизм площадки. Это место, где встречаются
  автор, автоматика и тот, кто решает. В описании пишут <<Closes \#42>> --
  и задача закроется сама, когда PR сольют.}


\begin{frame}{Actions}
  \slidetop
  \claim{Площадка \hi{сама запускает проверки} --- на каждое изменение, на
    чистой машине.}

  \begin{center}
    \begin{tikzpicture}[font=\small,
      st/.style={box,minimum width=2.6cm,minimum height=1.1cm}]
      \node[st] (a) {событие\\{\scriptsize\color{muted}открыт pull request}};
      \node[st,right=0.9cm of a] (b) {workflow\\{\scriptsize\color{muted}файл в репозитории}};
      \node[st,right=0.9cm of b] (c) {шаги\\{\scriptsize\color{muted}сборка, тесты}};
      \node[st,right=0.9cm of c,minimum width=2.2cm] (d)
        {\hi{прошло}\\\lo{не прошло}};
      \foreach \x/\y in {a/b,b/c,c/d} \draw[ar] (\x) -- (\y);
    \end{tikzpicture}
  \end{center}

  \slidegap
  \aside{Результат виден прямо в pull request. Собирает не автор и не на
    своём ноутбуке --- <<у меня работает>> здесь не проходит.}
\end{frame}
\note{GitHub Actions. Это и есть CI -- continuous integration. Сейчас
  достаточно принципа: событие, сценарий, вердикт. Как строить сам сценарий
  -- в разделе про защиту.}

\begin{frame}{Настройки и защиты}
  \slidetop
  \claim{Правила проекта --- это \hi{настройки репозитория}, а не
    договорённости в чате.}

  \begin{tabular}{@{}p{4.6cm}p{8.6cm}@{}}
    \textbf{\hi{роли}}            & кто читает, кто пишет, кто
      управляет \\[0.5em]
    \textbf{\hi{защита веток}}    & в main --- только через pull
      request \\[0.5em]
    \textbf{\hi{обязательные проверки}} & красный pull request слить
      нельзя \\[0.5em]
    \textbf{\hi{обязательное review}} & нужно одобрение ответственного \\[0.5em]
    \textbf{\hi{секреты}}         & ключи хранит площадка, а не код \\
  \end{tabular}

  \slidegap
  \aside{Правило, записанное в настройках, \hi{нельзя забыть} и нельзя
    обойти по дружбе.}
\end{frame}
\note{Settings: collaborators, rulesets, secrets. Показать вживую на своём
  репозитории. Теперь, когда понятно, где репозиторий живёт и что умеет
  площадка, -- что в нём должно лежать.}

\begin{frame}[fragile]{Анатомия репозитория}
  \slidetop
  \begin{minipage}[t]{0.40\textwidth}
\begin{lstlisting}[style=plain]
HelloWorld/
  .devcontainer/
  .github/
    workflows/
    CODEOWNERS
    pull_request_template.md
  docs/adr/
  src/
  tests/
  .gitignore
  CHANGELOG.md
  CONTRIBUTING.md
  LICENSE
  README.md
  packages.lock.json
\end{lstlisting}
  \end{minipage}\hfill
  \begin{minipage}[t]{0.55\textwidth}
    \vspace{0.3em}
    {\small
    \begin{tabular}{@{}p{2.9cm}p{4.3cm}@{}}
      \textbf{\hi{README}}        & что это и как запустить \\[0.45em]
      \textbf{\hi{CONTRIBUTING}}  & как сдавать работу \\[0.45em]
      \textbf{\hi{CODEOWNERS}}    & кто отвечает за какую часть \\[0.45em]
      \textbf{\hi{workflows}}     & автоматические проверки \\[0.45em]
      \textbf{\hi{docs/adr}}      & почему решили именно так \\[0.45em]
      \textbf{\hi{lock-файл}}     & точные версии зависимостей \\[0.45em]
      \textbf{\hi{CHANGELOG}}     & что менялось от версии к версии \\[0.45em]
      \textbf{\hi{.devcontainer}} & готовое окружение \\
    \end{tabular}}
  \end{minipage}
\end{frame}
\note{Кода здесь две папки из пятнадцати строк. Всё остальное -- то, что
  делает из папки с кодом продукт, в который можно войти. ADR --
  Architecture Decision Record: короткая запись <<что решили, какие были
  варианты, почему выбрали этот>>.}

\begin{frame}{<<Работает у меня>> --- не аргумент}
  \slidetop
  \claim{Сборка \hi{воспроизводима}: один коммит --- один результат, на
    любой машине.}

  \begin{tabular}{@{}p{4.6cm}p{8.6cm}@{}}
    \textbf{\hi{lock-файл}}      & зависимости те же, до последней
      цифры версии \\[0.55em]
    \textbf{\hi{версия SDK в репозитории}} & инструменты те же \\[0.55em]
    \textbf{\hi{dev-контейнер}}  & окружение то же: открыл --- и всё
      установлено \\[0.55em]
    \textbf{\hi{CI}}             & собирает не автор и не на его машине \\
  \end{tabular}

  \slidegap
  \aside{Проверка простая: чистая машина, \texttt{git clone}, одна команда.
    Не собралось --- \lo{это ошибка репозитория}, а не новичка.}
\end{frame}
\note{Dev-контейнер -- описание окружения файлом в репозитории: образ,
  версии, расширения редактора. Новый человек или агент получает рабочее
  место за минуты, а не за два дня переписки.}

\begin{frame}{Зависимости --- это чужой код в вашем продукте}
  \slidetop
  \claim{Вы отвечаете за всё, что собрали, --- включая то, что \lo{не
    писали}.}

  \begin{cbul}{accent}
    \item каждая зависимость \hi{добавляется осознанно}, отдельным pull
      request
    \item версии зафиксированы lock-файлом
    \item обновления приходят автоматически и проходят те же проверки
    \item известные уязвимости проверяются на каждой сборке
  \end{cbul}

  \slidegap
  \aside{ИИ может уверенно назвать пакет, которого \lo{не существует}.
    Имя такого пакета кто-нибудь зарегистрирует --- с чем угодно внутри.}
\end{frame}
\note{Dependabot и аналоги: бот сам открывает PR с обновлением. Выдуманные
  ИИ пакеты -- реальный вектор атаки: злоумышленник регистрирует имя,
  которое модели часто придумывают. Подробно о цепочке поставки -- в
  следующей лекции.}

\begin{frame}[plain]
  \keyline{После merge у кода нет автора. Это \hi{код репозитория}.}
\end{frame}
\note{Не вини автора кода. Сломалось -- значит, это пропустили проверки и
  review, то есть система. Чиним систему. Это мост к последнему разделу.}

% ================================================= 4. PROTECTION ===========
\section{4. Защита}

\begin{frame}[plain]
  \sectionopener{4}{Защита}{\hi{Never trust the contributor.}\\[0.3em]
    Ни новичку, ни эксперту, ни агенту, ни себе.}
\end{frame}
\note{Бюджет: 14 минут. Настройки мы уже видели; теперь -- какими они
  должны быть и почему. Это не про плохих людей. Ошибаются все; вопрос в
  том, где ошибка будет поймана. Вы должны быть жёсткими: дисциплина
  first.}

\begin{frame}{main всегда рабочий}
  \slidetop
  \claim{В любой момент main можно \hi{собрать и выпустить}. Это не
    договорённость, а настройка.}

  \begin{cbul}{accent}
    \item прямой push в main \lo{запрещён} --- всем, включая владельца
    \item изменения приходят только через pull request
    \item pull request нельзя слить, пока проверки не зелёные
    \item ветка перед слиянием \hi{актуальна}: проверен результат слияния,
      а не старая ветка
    \item нужно одобрение --- от человека, который отвечает за эту часть
  \end{cbul}
\end{frame}
\note{Branch protection, rulesets. Правило, которое можно обойти <<один раз,
  срочно>>, правилом не является. PR без задачи и без теста -- подозрительный
  PR.}

\begin{frame}{Fail fast}
  \slidetop
  \claim{Сначала дешёвые проверки. Упала первая --- \lo{дальше не идём}.}

  \begin{center}
    \begin{tikzpicture}[font=\small,
      st/.style={box,minimum width=2.0cm,minimum height=1.1cm}]
      \node[st,draw=accent,line width=1.0pt] (a) {\hi{формат}};
      \node[st,right=0.6cm of a] (b) {анализ};
      \node[st,right=0.6cm of b] (c) {сборка};
      \node[st,right=0.6cm of c] (d) {тесты};
      \node[st,right=0.6cm of d] (e) {безопасность};
      \foreach \x/\y in {a/b,b/c,c/d,d/e} \draw[ar] (\x) -- (\y);
      \node[muted,font=\scriptsize,below=0.5em] at (a.south) {секунды};
      \node[muted,font=\scriptsize,below=0.5em] at (e.south) {минуты};
    \end{tikzpicture}
  \end{center}

  \slidegap
  \hrulethin{0.42}
  \rulegap
  {\large Не отформатировано --- \hi{не заслуживает тестирования}.}
\end{frame}
\note{Зачем тратить десять минут тестов на код, который не прошёл проверку
  за пять секунд. И зачем человеку читать PR, который не прошёл автоматику.
  Статический анализ и сканирование безопасности -- часть разработки, а не
  отдельный этап <<потом>>.}

\begin{frame}{Где дешевле поймать ошибку}
  \slidetop
  \claim{Одна и та же ошибка стоит \hi{секунды} или \lo{деньги} --- смотря
    где её остановили.}

  \begin{center}
    \begin{tikzpicture}[font=\small,x=0.93cm,y=1cm]
      \foreach \i/\t/\w/\h in {0/pre-commit hook/у вас на машине/0.5,
                               1/CI/на сервере/1.0,
                               2/review/человек или ИИ/1.6,
                               3/пользователь/в продукте/2.6} {
        \fill[codebg] (\i*3.3,0) rectangle ++(2.7,\h);
        \draw[hairline,line width=0.8pt] (\i*3.3,0) rectangle ++(2.7,\h);
        \node[anchor=north,font=\small\bfseries,text height=1.6ex,text depth=0.5ex] at (\i*3.3+1.35,-0.1) {\t};
        \node[anchor=north,font=\scriptsize,text=muted] at (\i*3.3+1.35,-0.55)
          {\w};
      }
      \draw[accent,line width=1.0pt] (0,0) rectangle ++(2.7,0.5);
      \draw[alt,line width=1.0pt] (9.9,0) rectangle ++(2.7,2.6);
      \node[accent,font=\scriptsize\bfseries] at (1.35,0.25) {секунды};
      \node[font=\scriptsize] at (4.65,0.5) {минуты};
      \node[font=\scriptsize] at (7.95,0.8) {часы};
      \node[alt,font=\scriptsize\bfseries] at (11.25,1.3) {дни и деньги};
    \end{tikzpicture}
  \end{center}
\end{frame}
\note{Pre-commit hook -- скрипт, который Git запускает перед коммитом:
  формат, линтер, секреты. Но hook на машине можно отключить, поэтому он --
  удобство, а CI -- закон. Высота столбика условная: это порядок, а не
  измерение.}

\begin{frame}[fragile]{Репозиторий без CI --- свалка}
  \slidetop
  \begin{minipage}[t]{0.52\textwidth}
\begin{lstlisting}[style=plain,basicstyle=\ttfamily\scriptsize]
name: ci
on:
  pull_request:
    branches: [main]

jobs:
  build:
    runs-on: ubuntu-latest
    steps:
      - uses: actions/checkout@v4
      - uses: actions/setup-dotnet@v4
        with:
          global-json-file: global.json
      - uses: ./.github/actions/format
      - uses: ./.github/actions/build
      - uses: ./.github/actions/test
\end{lstlisting}
  \end{minipage}\hfill
  \begin{minipage}[t]{0.43\textwidth}
    \vspace{0.3em}
    \aside{Файл лежит в репозитории: проверки \hi{версионируются вместе с
      кодом}.}\par
    \vspace{0.9em}
    \aside{Шаги --- это \hi{actions}, а не россыпь команд: готовые и свои.}
      \par
    \vspace{0.9em}
    \aside{Своя action \hi{логически завершена и обособлена}: её можно
      взять в другой проект целиком.}\par
  \end{minipage}
\end{frame}
\note{GitHub Actions. Стройте CI из actions, а не из cli-команд в yaml:
  команду нельзя переиспользовать, протестировать и версионировать. Пишите
  свои actions. Три последние строки -- локальные actions этого же
  репозитория.}

\begin{frame}[plain]
  \keyline{Механизм не должен зависеть от доверия \hi{к конкретному
    человеку}.}
\end{frame}
\note{Итог лекции: репозиторий хранит всё, Git помнит всё, main защищён,
  ошибка ловится как можно раньше. В следующий раз -- тесты и review, а
  потом: кто всем этим управляет и как продукт доходит до пользователя.}

\begin{frame}[plain]
  \shout{Задавайте вопросы}
\end{frame}
\note{Если вопросов нет -- спросить зал: что из сегодняшнего вы можете
  включить в своём учебном репозитории сегодня вечером?}

\end{document}
